% Part:sets-functions-relations
% Chapter: size-of-sets
% Section: equinumerous-sets

\documentclass[../../../include/open-logic-section]{subfiles}

\begin{document}

\olfileid{sfr}{siz}{equ}

\olsection{సమసంఖ్యాకత}

పరిమిత సమితులకు చక్కగా పనిచేసే, సమితుల ``పరిమాణం'' గురించిన సహజమైన
భావన మనకు ఉంది. మరి అనంత సమితుల సంగతి ఏమిటి? \emph{ఎంత పరిమాణమైనా}
ఉండగల రెండు సమితుల పరిమాణాలను పోల్చడానికి ఒక ఆచారబద్ధమైన పద్ధతిని
రూపొందించాలంటే, సమితులు ఒకే పరిమాణంలో ఎప్పుడు ఉంటాయో నిర్వచించడం
మంచి ఆరంభం. ఫ్రేగె ఇలా అంటాడు:
\begin{quote}
  ఒక పరిచారకుడు బల్లపై పళ్లేల సంఖ్యకు కచ్చితంగా సమానమైన సంఖ్యలో
  కత్తులు పెట్టాడని నిర్ధారించుకోవాలనుకుంటే, ప్రతి పళ్లెం కుడివైపున
  ఒక కత్తిని, బల్లపైని ప్రతి కత్తీ ఏదో ఒక పళ్లెం కుడివైపున ఉండేలా
  అమర్చితే, వాటిలో దేనినీ లెక్కించనక్కరలేదు. అలా పళ్లేలు, కత్తులు
  ఒకదానితో మరొకటి ఏకైకంగా సహసంబంధించబడతాయి; అదీ అదే ప్రాదేశిక
  సంబంధం ద్వారా. \citep[\S70]{Frege1884}
\end{quote}
ఈ భాగంలోని అంతర్దృష్టిని ఒక ఆచారబద్ధ నిర్వచనం ద్వారా వ్యక్తం చేయవచ్చు:

\begin{defn}\ollabel{comparisondef}
  \tetoken{ఒక ద్విగుణ ప్రమేయం}{!!a{bijection}} $f \colon A \to B$ ఉంటేనే $A$, $B$తో
  \emph{సమసంఖ్యాకం} అంటాం; దీన్ని $\cardeq{A}{B}$ అని రాస్తాం.
\end{defn}

\begin{prop}\ollabel{equinumerosityisequi}
సమసంఖ్యాకత ఒక తుల్యతా సంబంధం.
\end{prop}

\begin{proof}
సమసంఖ్యాకత స్వావర్తన, సౌష్ఠవ, సంక్రామక ధర్మాలు గలదని చూపాలి.
$A, B$, $C$ సమితులు అనుకుందాం.

\emph{స్వావర్తనత్వం.} ప్రతి $x \in A$కు $\Id{A} (x) = x$ అయ్యే
తాదాత్మ్య చిత్రణ $\Id{A} \colon A \to A$ ఒక \tetoken{ద్విగుణ ప్రమేయం}{!!a{bijection}}.
కాబట్టి $\cardeq{A}{A}$.

\emph{సౌష్ఠవం.} $\cardeq{A}{B}$ అనుకుందాం; అంటే, ఒక
\tetoken{ద్విగుణ ప్రమేయం}{!!a{bijection}} $f\colon A \to B$ ఉంది. $f$ \tetoken{ద్విగుణ}{!!{bijective}} కనుక దాని
విలోమం $f^{-1}$ ఉంటుంది, అది కూడా \tetoken{ద్విగుణ}{!!{bijective}}. అందువల్ల
$f^{-1}\colon B \to A$ ఒక \tetoken{ద్విగుణ ప్రమేయం}{!!a{bijection}}; కాబట్టి $\cardeq{B}{A}$.

\emph{సంక్రామకత్వం.} $\cardeq{A}{B}$, $\cardeq{B}{C}$ అనుకుందాం;
అంటే, \tetoken{ద్విగుణ ప్రమేయాలు}{!!{bijection}s} $f\colon A \to B$, $g\colon B \to C$ ఉన్నాయి.
అప్పుడు సంయోగం $\comp{f}{g}\colon A \to C$ \tetoken{ద్విగుణ}{!!{bijective}}; కాబట్టి
$\cardeq{A}{C}$.
\end{proof}

\begin{prop}
$\cardeq{A}{B}$ అయితే, $A$ \tetoken{లెక్కించదగిన}{!!{enumerable}} కావడానికి అవసరమైన మరియు
సరిపడే షరతు $B$ లెక్కించదగినదిగా ఉండటం.
\end{prop}

\begin{editorial}
కింది నిరూపణ \olref[enm]{defn:enumerable}ను
\olref[enm]{sec} చేర్చబడితే, లేకపోతే
\olref[enm-alt]{defn:enumerable}ను వాడుతుంది.
\end{editorial}

\begin{proof}
$\cardeq{A}{B}$ అనుకుందాం; కాబట్టి ఏదో \tetoken{ద్విగుణ ప్రమేయం}{!!{bijection}}
$f \colon A \to B$ ఉంది. అలాగే $A$ \tetoken{లెక్కించదగిన}{!!{enumerable}} అనుకుందాం.
\oliflabeldef{sfr:siz:enm:defn:enumerable}{
  అప్పుడు $A = \emptyset$ లేదా \tetoken{ఒక సంగ్రస్త}{!!a{surjective}} ప్రమేయం
  $g\colon \PosInt \to A$ ఉంటుంది. $A = \emptyset$ అయితే
  $B = \emptyset$ కూడా (లేకపోతే \tetoken{ఒక మూలకం}{!!a{element}}~$y \in B$ ఉంటుంది,
  కానీ $f(x) = y$ అయ్యే $x \in A$ ఉండదు). మరోవైపు,
  $g\colon \PosInt \to A$ \tetoken{సంగ్రస్త}{!!{surjective}} అయితే,
  $\comp{g}{f} \colon \PosInt \to B$ \tetoken{సంగ్రస్త}{!!{surjective}}. దీన్ని చూడటానికి
  $y \in B$ తీసుకోండి. $f$ \tetoken{సంగ్రస్త}{!!{surjective}} కనుక $f(x) = y$ అయ్యే
  $x \in A$ ఒకటి ఉంటుంది. $g$ \tetoken{సంగ్రస్త}{!!{surjective}} కనుక $g(n) = x$ అయ్యే
  $n \in \PosInt$ ఒకటి ఉంటుంది. అందువల్ల,
\[
(\comp{g}{f})(n) = f(g(n)) = f(x) = y
\]
కాబట్టి $\comp{g}{f}$ \tetoken{సంగ్రస్త}{!!{surjective}}. $\comp{g}{f}$ అనేది $B$ యొక్క
లెక్కింపు; అందువల్ల $B$~\tetoken{లెక్కించదగిన}{!!{enumerable}}.}
{
  అప్పుడు $A = \emptyset$ లేదా, పరిధి $A$గానూ ప్రవేశం $\Nat$ లేదా
  సహజ సంఖ్యల ఒక ఆరంభ క్రమంగానూ ఉన్న \tetoken{ఒక ద్విగుణ ప్రమేయం}{!!a{bijection}}~$g$ ఉంటుంది.
  $A = \emptyset$ అయితే $B = \emptyset$ కూడా (లేకపోతే $f(x) = y$
  అయ్యే $x \in A$ ఏదీ లేని ఏదో $y \in B$ ఉంటుంది). కాబట్టి మనకు
  \tetoken{ద్విగుణ ప్రమేయం}{!!{bijection}}~$g$ ఉందనుకుందాం. అప్పుడు $\comp{g}{f}$ అనేది పరిధి
  $B$గానూ, ప్రవేశం $g$ ప్రవేశంతో సమానంగానూ (అంటే, $\Nat$ లేదా దాని
  ఒక ఆరంభ ఖండంగానూ) ఉన్న \tetoken{ఒక ద్విగుణ ప్రమేయం}{!!a{bijection}}. కాబట్టి $B$ \tetoken{లెక్కించదగిన}{!!{enumerable}}.}
\sourcecorrection{OLSIZ-006}{రెండు ఖాళీ-సమితి శాఖల్లోనూ ఆంగ్ల మూలం,
ఇంకా నిర్వచించని $g$ను వాడి $g(x)=y$ అని రాస్తుంది. అక్కడ అవసరమైన
ద్విగుణ ప్రమేయం $f$ కనుక, రెండు చోట్లా $f(x)=y$గా సరిచేశాం.}

$B$ \tetoken{లెక్కించదగిన}{!!{enumerable}} అయితే, $f$కు బదులు \tetoken{ద్విగుణ ప్రమేయం}{!!{bijection}}
$f^{-1}\colon B \to A$తో ఇదే వాదనను పునరావృతం చేసి $A$
\tetoken{లెక్కించదగిన}{!!{enumerable}} అని పొందుతాం.
\end{proof}

\begin{prob}
$\cardeq{A}{C}$, $\cardeq{B}{D}$, అలాగే $A \cap B =
C \cap D = \emptyset$ అయితే $\cardeq{A \cup B}{C \cup D}$ అని చూపండి.
\end{prob}

\begin{prob}
$A$ అనంతమైనది, \tetoken{లెక్కించదగిన}{!!{enumerable}} అయితే $\cardeq{A}{\Nat}$ అని చూపండి.
\end{prob}

\end{document}
